\chapter{Background and Related Work} \label{Chap2}\label{chap:background}

\section{Procedural and data-driven urban modelling}

Procedural modelling represents architecture through rules and parameters
rather than solely through sampled surfaces.  Early city systems used extended
L-systems to generate street networks and building populations
\cite{parish2001procedural}; split grammars then showed how attributes and
rules could control architectural variation at scale \cite{wonka2003instant}.
CGA shape grammar made the relationship between initial shapes, rule
operations, and polygonal buildings explicit \cite{muller2006procedural}.
ProceX made procedural extrusions
interactive: footprint-driven operations create architectural massing that can
be edited rather than merely viewed \cite{kelly2011procex}.  This structured
representation is attractive for planning, simulation, visualisation, and
content authoring because it preserves semantic parts.

Open building-footprint collections provide initial shapes for applying this
principle beyond manually authored scenes.  Microsoft's Global ML Building
Footprints dataset supplies predicted building outlines at large geographic
scale \cite{microsoftGlobalFootprints2026}.  These outlines provide location
and plan extent, while height, appearance, and architectural semantics must
still be reconstructed or generated.

Real data rarely arrives as clean procedural parameters.  BigSUR addresses a
closely related integration problem by combining noisy large-scale meshes,
imagery, and \gls{gis} footprints to recover structured urban models
\cite{kelly2017bigsur}.  FrankenGAN then enriches coarse mass models with
semantic fa\c{c}ade detail and appearance through style-synchronised generative
networks \cite{kelly2018frankengan}.  ChordAtlas exposes these research
components in an interactive Java environment, with generators for GIS data,
panoramas, meshes, and blocks \cite{kellyChordAtlas2024}.  FrankenGAN assumes a
structured building-mass representation: it is neither a topology-repair
method for arbitrary meshes nor a source of ground-truth fa\c{c}ade texture.
ChordAtlas appearance exports are therefore treated here as procedurally
generated assets rather than texture-preserving conversions of Sat3DGen
triangles.

Recent geospatial approaches divide geometry and appearance processing in
different ways.  GeoTexBuild begins with map footprints and separates
height-map generation, geometry reconstruction, and appearance stylisation
\cite{wang2025geotexbuild}.  Buyukdemircioglu and Oude Elberink instead texture
existing CityJSON geometry from oblique and nadir aerial imagery
\cite{buyukdemircioglu2024texture}.  Their contrasting starting points clarify
that geometric reconstruction, structured city representation, and appearance
transfer are related but distinct stages.

ArcGIS CityEngine is a relevant industrial authoring system in this procedural
lineage.  Its central workflow combines streets, blocks, lots, initial shapes,
GIS attributes, assets, and authored CGA rules to generate polygonal city
content \cite{esriCityEngineIntro2026}.  File-geodatabase attributes can be
retained and used to drive rules, while the documented runtime and SDK support
scripted generation and export \cite{esriCityEngineGIS2026,
esriCityEngineSDK2026}.  These capabilities favour controllable regeneration,
scenario authoring, and attribute-aware content production.  They do not,
however, imply automatic recovery of a real building from satellite pixels.
Likewise, importing a dense OBJ does not by itself create stable building
identity or roof, wall, and window semantics.  CityEngine therefore serves as
a representation and workflow comparator, including an executed local
example, but not as a satellite-reconstruction accuracy baseline.

This lineage establishes two requirements for the present integration.  First,
generated geometry must be spatially compatible with footprints and panoramas,
not merely visually plausible in isolation.  Secondly, it must be converted
into ChordAtlas's tiled \glspl{minimesh} and generator graph without corrupting
an existing project when conversion fails.  The baseline ChordAtlas workflow
can clip an already loaded mesh and construct a block.  The proposed path adds
external, on-demand generation and publication, so state and failure handling
become part of the research problem.

\section{Satellite-conditioned street-level geometry}

Satellite-to-ground synthesis is underconstrained: roofs and plan-view context
are visible from above, while fa\c{c}ades and street-level occlusions are not.
Cross-view datasets reduce that ambiguity by pairing overhead imagery with
ground panoramas.  VIGOR formalised a setting in which a query is not
necessarily centred in its reference satellite image and uses overlapping
reference tiles \cite{zhu2021vigor}.  Its correspondence geometry is therefore
not an incidental file format.  Candidate offsets must describe where the
ground query lies within each aerial tile, and sufficient overlap is required
if several candidates are meant to contain the same location.

Sat2Density learns a neural density representation from satellite--ground
pairs \cite{qian2023sat2density}.  Its successor extends the cross-view
rendering lineage \cite{qian2026sat2densitypp}.  Sat2Scene uses diffusion for urban scene
generation \cite{li2024sat2scene}, while Sat2City uses a cascaded latent
diffusion formulation for city generation \cite{hua2025sat2city}.  These
methods are related baselines, but their tasks, training data, and output
representations differ.  The comparison therefore considers their
architectures and evaluation protocols rather than treating their reported
metrics as interchangeable.

At the broader systems level, Meta's WorldGen and World Labs illustrate a move
towards prompt- and observation-conditioned 3D environments
\cite{metaWorldGen2025,worldLabs2026}.  Their relevance here is the objective of
producing navigable and editable world content from location-specific
observations; they provide systems context rather than directly comparable
Sat3DGen baselines.

Sat3DGen places particular emphasis on explicit geometry.  Its published
pipeline uses a frozen DINOv3 image encoder, spatial padding tokens, a decoder
that produces three feature planes, and a small network that predicts density
and colour \cite{simeoni2025dinov3,qian2026sat3dgen}.  The tri-plane design is
related to efficient geometry-aware neural rendering \cite{chan2022eg3d}.
Training adds a gravity-based density-variation loss, a relative-depth prior
derived from Depth Anything V2 \cite{yang2024depthanythingv2}, and randomly
projected perspective views from ground panoramas.  The density representation
can be rendered from satellite or street-level viewpoints and converted into
an explicit surface through Marching Cubes
\cite{qian2026sat3dgen,lorensen1987marching}.  Figure~\ref{fig:sat3dgen-model}
separates this published model from the large-image and geospatial extensions
studied here.

\begin{figure}[H]
\centering
\resizebox{0.98\textwidth}{!}{%
\begin{tikzpicture}[
  node distance=7mm and 6mm,
  model/.style={draw=evidenceblue,rounded corners,align=center,
    minimum height=11mm,text width=25mm,fill=evidenceblue!8},
  train/.style={draw=evidenceorange,rounded corners,align=center,
    minimum height=9mm,text width=25mm,fill=evidenceorange!8},
  downstream/.style={draw=evidencegreen,rounded corners,align=center,
    minimum height=11mm,text width=32mm,fill=evidencegreen!8},
  arrow/.style={-{Latex[length=2mm]},thick}
]
\node[model] (image) {Satellite RGB\\$256\!\times\!256$};
\node[model,right=of image] (dino) {Frozen DINOv3\\token encoder};
\node[model,right=of dino] (tokens) {$16\!\times\!16$ tokens\\+ spatial padding};
\node[model,right=of tokens] (planes) {Decoder\\three feature planes};
\node[model,right=of planes] (field) {MLP density\\and colour field};
\node[model,right=of field] (mesh) {Rendering or\\Marching Cubes};
\draw[arrow] (image) -- (dino);
\draw[arrow] (dino) -- (tokens);
\draw[arrow] (tokens) -- (planes);
\draw[arrow] (planes) -- (field);
\draw[arrow] (field) -- (mesh);
\node[train,below=11mm of dino] (gravity) {Gravity density\\variation loss};
\node[train,right=of gravity] (depth) {Relative-depth\\prior};
\node[train,right=of depth] (persp) {Perspective\\views from\\panoramas};
\draw[arrow,dashed,evidenceorange] (gravity.north) -- (tokens.south);
\draw[arrow,dashed,evidenceorange] (depth.north) -- (planes.south);
\draw[arrow,dashed,evidenceorange] (persp.north) -- (field.south);
\node[downstream,below=11mm of field] (large) {This dissertation:\\fractional density fusion\\and vertex colour};
\node[downstream,right=16mm of large] (bridge) {Geospatial mesh and\\ChordAtlas bridge};
\draw[arrow,evidencegreen] (field.south) --
  node[midway,right=1mm,font=\scriptsize,align=left]{overlapping\\windows}
  (large.north);
\draw[arrow,evidencegreen] (large.east) --
  node[midway,above=1mm,font=\scriptsize,align=center]{global\\coloured mesh}
  (bridge.west);
\draw[arrow,dashed,evidencegreen] (mesh.south) to[out=-90,in=90]
  node[pos=.55,right=1mm,font=\scriptsize,align=left]{independent\\mesh path}
  (bridge.north);
\end{tikzpicture}%
}
\caption{Published Sat3DGen components and the two project extensions.  Blue
and orange identify the model and training components of Qian et al.; green
identifies the large-image density/colour method and the subsequent geospatial
integration developed in this dissertation.  The dashed link is the separate
independent-mesh experiment \cite{qian2026sat3dgen}.}
\label{fig:sat3dgen-model}
\end{figure}

The paper reports training on VIGOR-derived cross-view pairs, evaluation
outside the training cities, and a large-scale strategy based on overlapping
patches.  Its analysis identifies pose accuracy, orthographic projection,
locally flat ground, and atypical architecture as important reconstruction
factors.  The present study extends the systems context by examining
local-to-world transforms, seam handling, DSM alignment, building association,
and ChordAtlas publication \cite{qian2026sat3dgen}.

Three processing domains must be distinguished.  The released Sat3DGen
large-image wrapper predicts overlapping density volumes, applies hard boundary
cropping and equal-weight averaging, and runs Marching Cubes once on the fused
field; unlike its single-image mesh path, the wrapper does not query or export
vertex colour \cite{qian2026sat3dgen,qianSat3DGenCode2026}.  The large-image
extension developed here retains one global density-domain surface extraction
but maps window locations at sub-voxel precision, feathers overlapping
contributions, and performs a memory-bounded second pass for vertex RGB.  The
subsequent geospatial path operates in mesh space, adding OSM/DSM alignment,
footprint association, per-building extraction, and ChordAtlas publication.

\section{Semantic city models, mesh validity, and provenance}

A renderable triangle soup is not yet a semantic city model.  CityGML 3.0
separates georeferenced geometry from urban-object semantics, relationships,
and multiple spatial representations \cite{ogcCityGML3}.  Its Levels of Detail
should not be reduced to a generic low/medium/high polygon-count switch:
geometric complexity and semantic granularity are distinct
\cite{biljecki2016lod}.  The bridge's footprint-linked identifiers are therefore a
useful coarse association; the current OBJ outputs do not encode the validated
roof and wall semantics required for a CityGML representation.

Downstream validity is similarly task-dependent.  Local mesh repair may close
holes or remove defects, but it can also change the observed surface
\cite{attene2010repair}; standards-oriented validation distinguishes valid
primitives from merely readable files \cite{ledoux2018val3dity}.  This motivates
reporting boundary and non-manifold edges separately from geometric accuracy.

Finally, a manifest is only a container unless it records causal lineage.
PROV-O distinguishes entities, activities, and agents \cite{w3cProvO2013}; in
this setting the equivalent minimum is an exact input allow-list and hashes,
geographic extent and CRS, code/model/configuration identity and execution state.  

\section{Geospatial contracts}

\subsection{Coordinate order and reference systems}

GeoJSON positions are written in longitude--latitude order
\cite{butler2016geojson}, whereas an Overpass QL bounding box is ordered south,
west, north, east \cite{overpassQL}.  Confusing the two conventions can send a
valid query to the wrong region without producing a type error.  This is a
representative integration defect: each subsystem may accept four floating
point values while disagreeing about their meaning.

The data-builder grid uses Web Mercator world pixels at zoom 20.  EPSG:3857 is
appropriate for web mapping, but its registry explicitly notes scale
distortion \cite{epsg3857}.  Metric analysis of the London DSM instead uses
EPSG:27700, the British National Grid \cite{epsg27700}.  The bridge establishes
a local frame whose origin is stored in WGS84 and whose axes are X east, Y
height, and Z south.  The frame is convenient for ChordAtlas and Sat3DGen mesh
operations, but a constant-metres-per-degree approximation is not equivalent
to a projected CRS.  Chapter~\ref{chap:results} quantifies this difference.

\subsection{Surface elevation and semantics}

The Environment Agency's \SI{1}{\metre} LiDAR Composite DSM records surface
height rather than bare-earth terrain.  It can therefore include buildings,
vegetation, and vehicles; its 2022 composite also joins surveys acquired at
different times \cite{environmentAgencyDSM2022}.  GeoTIFF provides the
georeferenced raster interchange contract \cite{ogc2019geotiff}.  A DSM can
help place generated geometry vertically, but it is not an exact and
contemporaneous building-height ground truth.

OSM supplies footprints and contextual categories.  It is a collaborative
database under the ODbL \cite{openstreetmapData}.  Overpass queries commonly
return ways, relations, and the supporting nodes needed to reconstruct their
geometry.  Those support nodes must not automatically become category features,
and multipolygon outer/inner rings must retain their roles.  The data audit in
Chapter~\ref{chap:results} shows why semantic masks cannot be accepted solely
because image files were produced.

\section{Synthesis and research gap}

The literature reviewed above addresses complementary parts of the target
workflow.  Cross-view models connect satellite imagery to density fields,
rendered views, or explicit meshes
\cite{qian2023sat2density,qian2026sat3dgen}; procedural systems support
structured and editable urban content
\cite{kelly2017bigsur,kellyChordAtlas2024,esriCityEngineIntro2026}; and OSM,
DSM, and city-model standards provide geographic context, elevation, and
semantic conventions
\cite{openstreetmapData,environmentAgencyDSM2022,ogcCityGML3}.  These
components use different spatial frames and representations, so their outputs
do not compose automatically.

Among the reviewed approaches, none follows overlapping
satellite-conditioned predictions through regional fusion, geospatial
registration, footprint-based building extraction, and publication into an
interactive procedural environment.  Neural reconstruction work primarily
evaluates geometry or rendering, whereas procedural workflows generally begin
with geometry and geographic structure already available.  The research gap
is therefore the transition between these representations and the evaluation
of that transition as an integrated workflow.

Accordingly, this study evaluates whether satellite tiles, projected DSM data,
local mesh coordinates, and ChordAtlas placement remain spatially consistent.
It measures the effects of fusion, cropping, stitching, vertical correction,
and building extraction.  These evaluation
dimensions are operationalised in Chapter~\ref{chap:methodology} and provide
the basis for the research question introduced in
Chapter~\ref{chap:introduction}.
